\subsection{Strict positivity from the tautological metric}
\label{subsec:tautological-homogeneous-positivity}
\index{Vector bundle!tautological positivity}
\index{Tautological line bundle!homogeneous Levi form}

The second part of \cref{cor:continuous-convex-norm-obstruction} requires a smooth strictly plurisubharmonic function on $F^\times$ that is complex homogeneous of degree two. Such a function is the squared norm of a vector, viewed as a point of the tautological line it spans, for a suitable metric on the tautological line bundle. Its Levi form is strictly positive exactly when the dual tautological metric is positively curved, so the function exists when $F^*$ is ample.

Let $F\to X$ be a holomorphic vector bundle of positive rank. In this subsection, $\mathbb P(F)$ parametrizes complex lines in $F$, $\mathcal O_{\mathbb P(F)}(-1)$ is the tautological line bundle whose fiber over a line is that line, and
\[
 L=\mathcal O_{\mathbb P(F)}(1)
   =\mathcal O_{\mathbb P(F)}(-1)^*.
\]
The map $v\mapsto([v],v)$ identifies $F^\times$ holomorphically with the complement of the zero section in $\mathcal O_{\mathbb P(F)}(-1)$. We use the curvature convention $R_h=-\ddc\log h$ of \cref{sec:convex-bundle-norm-obstruction}.

\Needspace{5\baselineskip}
\begin{lemma}\label{lem:tautological-homogeneous-levi}
A smooth Hermitian metric $h$ on $L$ determines a positive smooth complex homogeneous function $H$ of degree two on $F^\times$: the squared norm of $v$ in the dual metric on the tautological line $\mathcal O_{\mathbb P(F)}(-1)_{[v]}$. The function $H$ is strictly plurisubharmonic if and only if $R_h$ is strictly positive on $\mathbb P(F)$.
\end{lemma}
\begin{proof}
Take local holomorphic coordinates $y$ on $\mathbb P(F)$ and a holomorphic frame $e$ of the tautological line bundle, and write the metric of the dual frame as $h=\mathrm e^{-\varphi}$. The dual metric gives
\[
 H(y,\lambda e)=|\lambda|^2\mathrm e^{\varphi(y)},\qquad
 R_h=\ddc\varphi.
\]
Let $\mathcal L$ denote the complex Hessian, without the factor $\mathrm i/(2\pi)$ of the curvature normalization. Differentiating at a point with $\lambda\ne0$ gives
\[
 \mathcal L H(u,t)
 =H\left(\mathcal L\varphi(u)+
          \left|\frac{t}{\lambda}+\partial\varphi(u)\right|^2\right),
 \qquad (u,t)\in T_y^{1,0}\mathbb P(F)\oplus\mathbb C.
\]
If $\mathcal L\varphi$ is positive definite, the right-hand side is positive for every nonzero $(u,t)$: the first term detects $u\ne0$, and the second detects $t\ne0$ when $u=0$. Conversely, for $u\ne0$, set $t=-\lambda\partial\varphi(u)$; strict positivity of $\mathcal L H$ then forces $\mathcal L\varphi(u)>0$. This is exactly positivity of $R_h$. The squared tautological norm is independent of the chosen frame and is complex homogeneous, so the local functions define the asserted global function $H$.
\end{proof}

A line bundle is very ample when its complete linear system defines a holomorphic embedding, and ample when some positive tensor power is very ample. With the convention above, a vector bundle $F^*$ is called ample when $L$ is ample.

\Needspace{5\baselineskip}
\begin{proposition}\label{prop:ample-dual-strict-homogeneous-function}
\index{Ampleness!strict homogeneous functions on the dual}
If $F^*$ is ample, there is a positive complex homogeneous function $H$ of degree two on $F^\times$ that is smooth and strictly plurisubharmonic there and, after extension by zero, continuous on $F$.
\end{proposition}
\begin{proof}
Choose $m\geq1$ such that $L^m$ is very ample. Its complete linear system gives an embedding $\iota\colon\mathbb P(F)\to\mathbb P^N$ with $L^m=\iota^*\mathcal O_{\mathbb P^N}(1)$. Pull back the Fubini--Study metric to a metric $k$ on $L^m$, and take its positive $m$th root $h$ on $L$: in a local frame $h=k^{1/m}$, and the transformation law for $L^m$ makes these local roots a metric on $L$. Its curvature is
\[
 R_h=\frac1m\iota^*\omega_{\mathrm{FS}}>0,
\]
where $\omega_{\mathrm{FS}}$ is the curvature form of the Fubini--Study metric on $\mathcal O_{\mathbb P^N}(1)$; the pullback is strictly positive because the differential of the embedding is injective. Apply \cref{lem:tautological-homogeneous-levi}. Since the base is compact, the resulting $H$ is bounded on the unit sphere bundle of a fixed Hermitian metric. Homogeneity then gives $0<H(v)\leq C|v|^2$ off zero, which proves continuity of the extension by zero.
\end{proof}

\Needspace{5\baselineskip}
\begin{corollary}\label{cor:ample-dual-continuous-convex-obstruction}
Let $F$ carry a holomorphic structure in the neighborhood of \cref{thm:convex-norm-deformation-obstruction}. If $F^*$ is ample, there is no positive continuous complex homogeneous function of degree two on $F$ that is plurisubharmonic off zero and locally strongly convex in the fibers.
\end{corollary}
\begin{proof}
By \cref{prop:ample-dual-strict-homogeneous-function}, ampleness of $F^*$ gives the function $H$ required in \cref{prop:weak-homogeneous-convex-replacement}, so the second assertion of \cref{cor:continuous-convex-norm-obstruction} applies.
\end{proof}

The function $H$ is not the square of a Hermitian norm in general, since the metric on the tautological line may vary with the line. It provides the strict plurisubharmonicity on the total space needed for regularization, whereas local strong convexity in the fibers remains a separate hypothesis in the continuous obstruction.
